\subsection{再论匹配圈}\label{sec:18k}
取\stref{定理}{18.24}中拟等价的一个逆拟等价，将其记为
\[
\G:\Tw\F(\pi)\longrightarrow\Tw\F^\to(\symbfit V_{\symbfit{\gamma}}^\#).
\]
首先注意，可以确定每个 (B) 型等变膜在 $\G$ 下映到哪里。由构造，对用于定义 $\G$ 的那个特定基，有
\[
\G(\widetilde\Delta_{\gamma_j}^\#)=V_{\gamma_j}^\#,
\]
其中右端是作为 $\F^\to(\symbfit V_{\symbfit{\gamma}}^\#)$ 本身的第 $j$ 个对象的消失圈。现考虑任意对象集 $\symbfit{\delta}$，及其第 $k$ 次初等 Hurwitz 移动 $\symbfit{\rho}=\mathcal L_k\symbfit{\delta}$。利用\stref{引理}{5.6}及\stref{命题}{18.23}，得到
\begin{equation}\label{eq:18.30}
\G(\widetilde\Delta_{\rho_j}^\#)\cong
\begin{cases}
\G(\widetilde\Delta_{\delta_j}^\#),&j\ne k,k+1,\\
\G(\widetilde\Delta_{\delta_k}^\#),&j=k+1,\\
T_{\G(\widetilde\Delta_{\delta_k}^\#)}\G(\widetilde\Delta_{\delta_{k+1}}^\#),&j=k.
\end{cases}
\end{equation}
这里的 $\cong$ 在 $D\F^\to(\symbfit V_{\symbfit{\gamma}}^\#)$ 中成立，而且实际可能只在相差平移的意义下成立，除非额外追踪正确的分次；我们不坚持这样做。\znote{289}{\eref{18.30} 最容易看反的是 $j=k$ 与 $j=k+1$ 两行的分工：初等 Hurwitz 移动把第 $k$ 条与第 $k+1$ 条路径换位，于是新基的第 $k+1$ 个对象就是原来的 $\widetilde\Delta_{\delta_k}^\#$（不加扭转），只有第 $k$ 个对象被扭转。逆移动公式中出现 $T^\vee$，方向相反，正说明 $\mathcal L_k$ 与 $\mathcal L_k^{-1}$ 互逆。另需留意扭转的下标是像对象 $\G(\widetilde\Delta_{\delta_k}^\#)$ 而不是原对象，因为变异是在目标范畴中进行的。「只到相差平移」并不削弱结论——后续比较都允许对象平移——但若把它当作严格同构去复合，分次上必出错。[来源：原书 §18k（印刷页 288、289）]}通过代入，可得到逆 Hurwitz 移动 $\symbfit{\sigma}=\mathcal L_k^{-1}\symbfit{\delta}$ 的类似公式，即
\[
\G(\widetilde\Delta_{\sigma_j}^\#)\cong
\begin{cases}
\G(\widetilde\Delta_{\delta_j}^\#),&j\ne k,k+1,\\
\G(\widetilde\Delta_{\delta_{k+1}}^\#),&j=k,\\
T^\vee_{\G(\widetilde\Delta_{\delta_{k+1}}^\#)}\bigl(\G(\widetilde\Delta_{\delta_k}^\#)\bigr),&j=k+1.
\end{cases}
\]
每条消失路径都是某个特选基的一部分，这就完成了第一步的说明。

\SourcePage{297}
其次，设 $\mu$ 为匹配路径，它相对于 $\symbfit{\delta}$ 的位置如\sr{fig:18.3}{图~18.3} 所示。由\stref{引理}{18.18}，$D\F(\pi)$ 中从 $\widetilde\Delta_{\delta_k}^\#$ 到 $\widetilde\Delta_{\delta_{k+1}}^\#$ 的态射空间有两个生成元，次数相差 $n$。准确的绝对次数
\SourceFigure{figure-18-3}{0.191076717}{18.3}
取决于分次的选择；我们平移这些分次，直到两个次数成为 $0$ 与 $n$。取匹配圈 $\Sigma_\mu$ 及相应的 (U) 型等变膜 $\widetilde\Sigma_\mu^\#$。由\stref{推论}{3.9}与\stref{命题}{18.21}，有
\[
\G(\widetilde\Sigma_\mu^\#)
\cong\G\bigl(\Cone(\widetilde\Delta_{\delta_k}^\#\to\widetilde\Delta_{\delta_{k+1}}^\#)\bigr)
\cong\Cone\bigl(\G(\widetilde\Delta_{\delta_k}^\#)\to\G(\widetilde\Delta_{\delta_{k+1}}^\#)\bigr),
\]
其中锥在任意非平凡零次态射上构造，同构的意义与\eref{18.30}相同。注意，对任意给定匹配路径 $\mu$，都能找到使这一做法适用的特选基 $\symbfit{\delta}$，这意味着全部匹配圈在 $\G$ 下的像都可显式确定。应用\stref{命题}{18.13}，得到：

\StatementNumber{26}\begin{corollary}\label{st:18.26}
设 $\mu_1,\ldots,\mu_l$ 为 $\pi$ 的有限条匹配路径。那么，只需在拟同构意义下知道一个特选基 $\symbfit{\gamma}$ 对应的 $\F^\to(\symbfit V_{\symbfit{\gamma}}^\#)$，就能计算 $\F(E)$ 中由匹配圈 $\Sigma_{\mu_j}^\#$ 构成的全 $\Ai$ 子范畴，计算结果确定到拟同构以及可能的对象平移。
\end{corollary}
